\documentclass[11pt,letterpaper]{article}
\usepackage[lmargin=1in,rmargin=1in,bmargin=1in,tmargin=1in]{geometry}
\usepackage{checkins}


% -------------------
% Content
% -------------------
\begin{document}
\thispagestyle{title}

% 08/24
\checkin{08/24} If $f(x)$ is a relation with $f(1)= 8$, $f(2)= 8$, and $f(3)= 3$, then $f(x)$ cannot be a function. \pspace

\sol The statement is \textit{false}. Recall a function is a relation such that for each input, there is only one possible output (though that output can be repeated by other inputs). That is, we can think of a function as a `coherent set of instructions.' Given an input, is there a single, unambiguous output? We know $f(1)= 8$, $f(2)= 8$, and $f(3)= 3$. There is no confusion as to what the output for $f(x)$ is for $x= 1, 2, 3$. In fact, $f(x)= -\tfrac{5}{2} x^2 + \tfrac{15}{2} x + 3$ is a function with $f(1)= 8$, $f(2)= 8$, and $f(3)= 3$. \par\vspace{1.3cm}



% 08/26
\checkin{08/26} The function $\ell(x)= 9 \;-\; 4x$ is linear, has $y$-intercept $9$, has slope $4$, and is an increasing function. \pspace

\sol The statement is \textit{false}. We know a linear function has the form $y= mx + b$. Observe that $\ell(x)$ has the form $mx + b$ with $m= -4$ and $b= 9$. Therefore, $\ell(x)$ is linear. We know the slope is $m$, and $m= -4$ (not $4$), and the $y$-intercept is $b$, and $b= 9$. We also know the $y$-intercept occurs when the input is $0$, i.e. $x= 0$. But we know $\ell(0)= 9 - 4(0)= 9 - 0= 9$. Finally, we know a linear function is increasing if $m > 0$, constant if $m= 0$, and decreasing if $m < 0$. Because $\ell(x)$ has $m= -4 < 0$, we know that $\ell(x)$ is decreasing. \par\vspace{1.3cm}



% 08/31
\checkin{08/31} A curve which `bends upwards' is convex. \pspace

\sol The statement is \textit{true}. A curve which `bends upwards' is said to be concave up or convex. A curve which `bends downwards' is concave down or concave. However, this is \textit{not} the same as increasing/decreasing. A curve can be either increasing/decreasing while also being either convex/concave. \par\vspace{1.3cm}



% 09/02
\checkin{09/02} If $W(a)= 11.5a - 10$ represents a linear model to predict the average child of age $a$'s weight in a local community, then the slope represents the weight gained each year and the $y$-intercept represents the initial weight. \pspace

\sol The statement is \textit{false}. It is generally true that a $y$-intercept represents an initial value. This is because the $y$-intercept occurs when the independent (input) variable is $0$. This typically will represent some type of `initial state output', as the input variable often represents something akin to time. However, this may not always be the case. Furthermore,  even if this is the case, that does not mean that the value has a sensible interpretation in context. In this instance, the $y$-intercept occurs when $a= 0$, i.e. when the age is $0$. However, how exactly can a child be aged $0$? At best, this is because they have just been born. But in this case, the model then predicts they have weight $W(0)= 11.5(0) - 10= 0 - 10= -10$~lbs. However, a person cannot have negative weight. So, even in this case, the value has no interpretation in context. This does not necessarily make the model poor, however. The domain is likely restricted to values where the model outputs good values. Notice that for values of $a \in [8, 15]$, the model outputs reasonable weights for children of that age. \par\vspace{1.3cm}



\newpage
























\end{document}